% =============================================================
% Paper Template (LaTeX) – "Derivations in Historical Context"
% Author: Christian Weilharter
% Series: Theoretical Physics – Derivations
% Purpose: 6×9 in (KDP/Print) + EPUB-compatible PDF (Apple Books)
% =============================================================
\documentclass[12pt,oneside]{scrartcl}
\usepackage[utf8]{inputenc}
\usepackage[T1]{fontenc}
\usepackage[english]{babel}
\usepackage{../styles/paper-style-en}

% ---- Literatur (BibLaTeX/Biber) ----
\usepackage[backend=biber,style=numeric,sorting=nyt]{biblatex}
\renewcommand*{\bibfont}{\small}
\addbibresource{../bib/literatur.bib}
% =============================================================
% Series and publication metadata
% =============================================================
\newcommand{\SeriesTitle}{Theoretical Physics – Derivations in Historical Context}
\newcommand{\PaperTitle}{The Compton Effect – Derivation, History, and Applications}
\newcommand{\PaperAuthor}{Christian Weilharter}
\newcommand{\PaperDate}{October 26, 2025}

\newcommand{\PaperDOI}{\emph{DOI (Zenodo)}: \texttt{10.5281/zenodo.xxxxxxx}}
\newcommand{\PaperISBNPDF}{\emph{ISBN (Print)}: 978-3-xxxx-xxxx-x}
\newcommand{\PaperISBNEPUB}{\emph{ISBN (E-Book)}: 978-3-xxxx-xxxx-x}
\newcommand{\PaperKeywords}{Compton Effect; Quantum Physics; X-ray Scattering; Photon Momentum; Four-Vectors}

% ---------- Metadata block ----------
\newcommand{\PaperMetadata}{
	\vspace{1em}
	\begin{flushleft}
		\small
		\setlength{\parskip}{3pt}
		\noindent\PaperDOI\\
		\PaperISBNPDF\\
		\PaperISBNEPUB\\
		\textit{Keywords:} \PaperKeywords
	\end{flushleft}
	\vspace{1em}
}

% =============================================================
% Title
% =============================================================
\title{\PaperTitle\\[0.5em]\large \SeriesTitle}
\author{\PaperAuthor}
\date{\PaperDate}

\begin{document}
	\thispagestyle{empty}
	\begin{center}
		{\Huge\bfseries \PaperTitle\par}
		\vspace{0.5em}
		{\Large\itshape \SeriesTitle\par}
		\vspace{2em}
		{\Large \PaperAuthor\par}
		\vspace{0.5em}
		{\large \PaperDate\par}
	\end{center}
	
	\PaperMetadata
	
	\vspace{1em}
	\hrule
	\vspace{1em}
	
	\noindent\textbf{Abstract.}  
	This paper combines a compact derivation of the Compton formula with the historical context from Planck to Compton and outlines its applications in modern physics. The focus is on clarity, mathematical precision, and the integration of physical theory into its historical development.
	
	\vspace{1em}
	\hrule
	\vspace{2em}
	
	% =============================================================
	% Main content
	% =============================================================
	
	
	
	\section{Introduction}
	
	When Arthur Compton investigated the scattering of X-rays by electrons in 1923, he encountered a result that the classical wave theory of light could not explain. The wavelength of the scattered radiation was greater than that of the incident beam — and the difference depended on the scattering angle. This experiment provided the first direct evidence for the transfer of momentum between photon and electron, thereby confirming the particle nature of light.
	
	The Compton effect marks a turning point in physics: it links energy and momentum conservation on a quantum-mechanical level and leads to the well-known \textbf{Compton formula}, which precisely describes the wavelength shift. The following derivation develops this relationship step by step — historically framed, mathematically sound, and physically transparent.
	
	\vspace{1.2em}
	\begin{quote}
		\small
		\textbf{Note on Methodology.}  
		This paper was created in close collaboration between human and artificial intelligence. 
		The physical reasoning and overall structure originate entirely from the author; the AI served as a tool for formulation, refinement, and didactic clarity. 
		The aim is to demonstrate that modern AI assistance does not replace human scientific work but enhances it — through greater clarity, structure, and transparency.
	\end{quote}
	\vspace{1.2em}
	
	
	\section{Historical Background}
	A \cite{einstein1909}concise timeline (2–3 pages): Planck (1900), Einstein (1905), Millikan (1916/17), Compton (1923). For each stage: the problem, the idea, and the impact. Optionally: brief quotes or an illustration of the experimental setup.
	\begin{PhysicsBox}[Photon–Electron Interaction]
		The Compton effect shows that photons transfer momentum to electrons,
		proving the particle aspect of light.
	\end{PhysicsBox}
	
	\begin{MathBox}[Compton Formula]
		\[
		\Delta\lambda = \frac{h}{m_e c}\,(1 - \cos\theta)
		\]
	\end{MathBox}
	\section{Physical Foundations}
	Photon momentum, conservation of energy and momentum, scattering geometry (definition of the angle $\theta$). Brief derivation of $p = \tfrac{h}{\lambda}$ and its connection to energy $E = hc/\lambda$.
	
	\begin{figure}[h]
		\centering
		\fbox{\parbox{0.85\linewidth}{\centering Figure: Geometry of Compton scattering (placeholder).}}
		\caption{Incident photon, scattering angle $\theta$, and outgoing electron direction.}
		\label{fig:geometry}
	\end{figure}
	
	\section{Derivation of the Compton Formula}
	Energy and momentum conservation before and after the collision, elimination of the electron velocity, and conversion to wavelength form:
	\begin{equation}
		\Delta\lambda = \lambda' - \lambda = \frac{h}{m_e c}\,(1 - \cos\theta),
		\label{eq:compton}
	\end{equation}
	where $\lambda_C = \tfrac{h}{m_e c}$ is the Compton wavelength.
	
	\section{Example Calculation}
	Choose, for instance, $\lambda = \SI{0.071}{\nano\meter}$ and $\theta = 90^\circ$. Substitute into \cref{eq:compton} and discuss the order of magnitude of $\Delta\lambda$.
	
	\section{Discussion and Limits}
	Thomson limit ($\lambda \gg \lambda_C$), relativistic corrections at high energies, experimental confirmations, and typical sources of error.
	
	\section{Applications}
	X-ray scattering, material analysis (e.g., electron density determination), astrophysical phenomena.
	
	\section{Conclusion}
	A concise summary (about 150–200 words): What does the Compton effect reveal about the nature of light? Relevance to modern physics.
	
	\section*{Acknowledgments (optional)}
	Brief thanks to collaborators, institutions, or funding sources.
	
	% =============================================================
	% References
	% =============================================================
%\section*{Literatur}
%\section*{References}
	\renewcommand*{\bibfont}{\small}
\printbibliography

	
	% ==========================================
	% Back page (final page) – revised, without repetition
	% ==========================================
	\clearpage
	\thispagestyle{empty}
	
	\vspace*{1.5cm}
	
	{\Large\bfseries The Compton Effect}\\[0.5em]
	{\large\itshape Derivation, History, and Applications}
	
	\vspace{1.2cm}
	
	\noindent
	\textbf{About this Paper}\\
	The Compton effect marks a turning point in physics:
	For the first time, it was shown unambiguously that light not only
	exhibits wave properties but also transfers momentum—like a particle.
	This paper presents a clear and concise derivation of the Compton
	formula, complemented by its historical background and significance
	for modern physics.
	
	\vspace{1.0cm}
	
	\noindent
	\textbf{Open Access}\\
	The \emph{free} PDF version is available at:\\[2pt]
	\href{https://doi.org/10.5281/zenodo.xxxxxxx}{\texttt{https://doi.org/10.5281/zenodo.xxxxxxx}}
	
	\vspace{0.8cm}
	
	\noindent
	\textbf{Series}\\
	\textit{Theoretical Physics – Derivations in Historical Context}\\
	More papers and materials:\\
	\href{https://mathandphysics.eu}{\texttt{mathandphysics.eu}}
	
	\vfill
	
	\noindent
	\small
	\textit{Bibliographic Information}\\
	\PaperDOI\\
	\PaperISBNPDF\\
	\PaperISBNEPUB
	
	\vspace{0.8cm}
	
	\noindent
	\footnotesize
	© 2025 Christian Weilharter — All rights reserved.\\
	Printed on demand via Amazon KDP. Typeset in \LaTeX.
	
	\vfill
	\begin{center}
		\small
		\textit{Math \& Physics Paper Series – Open Access at}\\[2pt]
		\href{https://mathandphysics.eu}{\texttt{https://mathandphysics.eu}} \, | \,
		\href{https://zenodo.org/communities/mathandphysics}{\texttt{Zenodo: Math \& Physics Open Science Collection}}
	\end{center}
	
\end{document}
